\problem{Жұлдыз тұрақтылығы}{10.0}

\pic{9}{R}{0pt}{0.263}{figures/2.1.png}

Жұлдыз --- түнгі аспандағы мыңдаған жарық нүктелердің бірі ғана емес, сонымен қатар астрофизикалық зерттеулер үшін маңызды нысан. Бұл есеп жұлдыздың түзілуі мен тепе-теңдігін сақтаудың негізгі механизмдерін зерттеуге бағытталған. Гравитациялық тұрақты $G=\qty{6.674e-11}{\Ggrav}$, жарық жылдамдығы $c=\qty{3.00e8}{\m\per\s}$, Больцман тұрақтысы $k=\qty{1.38e-23}{\J\per\K}$, Стефан--Больцман тұрақтысы $\sigma=\qty{5.67e-8}{\W\per\m\squared\per\K\tothe4}$, атомдық масса бірлігі $\qty{1}{\amu}=\qty{1.66e-27}{\kg}$.

Астрофизикада ``$\odot$'' индексі берілген шаманың біздің Күнге қатысты екенін білдіреді. Бұл есепте келесі мәліметтерді белгілі деп есептеңіз: Күннің массасы $M_\odot=\qty{2e30}{\m}$, Күннің радиусы $R_\odot=\qty{6.96e8}{\m}$, сондай-ақ Жердің орташа орбиталық радиусы $R_E=\qty{1.5e11}{\m}$.

\part{Параллакс}

\pic{10}{R}{0pt}{0.24}{figures/2.2.png}

Жұлдыздық параллакс --- жұлдыздың біздің Күн жүйесінен қашықтығын анықтаудың ең алғашқы әдісі. Мысалы, 1989 жылы астрометриялық деректерді жинау үшін ұшырылған Hipparcos (\textbf{Hi}gh \textbf{P}recision \textbf{Par}allax \textbf{Co}llecting \textbf{S}atellite) серігі параллакстарды миллиарксекундқа дейінгі дәлдікпен анықтай алады. Оң жақтағы суретте көрсетілгендей, Жердің орбиталық жазықтығына перпендикуляр орналасқан $S$ жұлдызын қарастырайық (мұндай жұлдызға, мысалы, $\delta$ Дракона жатады). Оны бір жыл бойы бақылағанда, ол қозғалмайтын алыс жұлдыздарға қатысты $C$ шеңбері бойымен ``қозғалады''. Hipparcos деректеріне сәйкес, $C$ орбитасының $\Delta\varphi$ бұрыштық диаметрі $\num{66.7e-3}$ бұрыштық секундқа тең.

\begin{task}
\item Күннен $S$ жұлдызына дейінгі $d$ қашықтықты жарық жылымен анықтаңыз.
\end{task}

\part{Протожұлдыздың түзілуі}

\pic{10}{R}{0pt}{0.325}{figures/2.3.png}[W51 тұмандығы --- Құс жолындағы ең үлкен жұлдыз фабрикаларының бірі]

Жұлдыздардың көпшілігі жұлдызаралық кеңістіктегі өте суық (шамамен 10–30 К) молекулалық бұлттардың арасында кластерлер болып түзіледі. Гравитациялық тартылыс әсерінен газ бұлты сығыла бастап, протожұлдызды құрайды. Радиусы $R$ (оны \taskref{2}–\taskref{4} пункттерінде белгілі деп санауға болады) және орташа температурасы $T=\qty{10}{\K}$ (мұндай температурада молекулалардың айналмалы еркіндік дәрежелері ``қатып қалған'') біртекті шар түріндегі сутегі бұлтының жеңілдетілген моделін қарастырайық. Шардың әрбір молекуласының орташа массасы $m=\qty{2}{\amu}$, олардың концентрациясы $n=\qty{4e8}{\m\tothe{-3}}$. Шын мәнінде, протожұлдыздың түзілу механизмі әлдеқайда күрделі: коллапсқа газ қысымы, магнит өрісінің болуы, центрифугадағы күштер мен турбуленттілік кедергі келтіреді, сондықтан коллапс әдетте соққы толқындарының таралуынан басталады.

\begin{task}
\item Бұлт молекулаларының жылулық қозғалысының толық $U$ ішкі энергиясын өрнектеңіз.
\item Бұлт ішіндегі еркін түсу үдеуі $g(r)$ радиус $r$-ге және оның толық $M$ массасына қалай тәуелді екенін анықтаңыз.
\item Массасы $M$ және радиусы $R$ шардың меншікті гравитациялық энергиясы $E_G$ келесі түрде беріледі:
$$
E_G = -f\frac{GM^2}{R},
$$
мұндағы $f$ --- өлшемсіз коэффициент, жалпы жағдайда шар ішіндегі массаның таралуын көрсетеді. Біртекті шар жағдайында $f$-ті анықтаңыз.
\item Бұлт коллапсының міндетті шарты --- толық $E$ механикалық энергиясының теріс болуы. Осыдан берілген сандық параметрлер бойынша бұлттың $M_J$ критикалық толық массасын $M_\odot$ бірлігімен анықтаңыз. Бұл шама Джинс массасы деп аталады --- $J$ индексі содан шыққан.
\end{task}

\part{Күн параметрлері}

Назарымызды Күнге аударайық. Басқа жұлдыздар сияқты, ол өз ішіндегі $P(r)$ қысым градиентінің болуына байланысты өз гравитациялық тартылысынан коллапсқа ұшырамайды. Бұл бөлімде Күн радиациясының $P(r)$-ге қосатын үлесін, сондай-ақ конвективті эффектілерді ескермеңіз. \textbf{Назар аударыңыз:} \taskref{2}–\taskref{5} пункттерінің нәтижелері бұдан былай жарамсыз, өйткені Күн ішіндегі заттың таралуын біртекті деп санауға \textbf{болмайды}!

\begin{task}
\item $r<R_\odot$ қашықтықтағы тығыздығы $\rho(r)$ және қалыңдығы $\Delta r$ болатын Күн элементін диск түрінде қарастырып, гидростатикалық тепе-теңдікті талдаңыз. Осыдан радиалық қысым градиенті $dP/dr$-дің $r$, $\rho(r)$ және радиусы $r$ болатын сфералық қабықшаның ішіндегі Күн массасының бөлігі $M(r)$-ге тәуелділігін қорытып шығарыңыз (осылайша, $M(R_\odot)=M_\odot$).
\end{task}
Қысым мен тығыздық қашықтықтың күрделі функциялары болғандықтан, Күннің көлемі бойынша орташаланған $\langle P\rangle$ қысымын қолдану арқылы есептеулерді жеңілдетуге болады:
$$
\langle P\rangle = \frac1{V_\odot} \int_0^{V_{\odot}} P\, dV.
$$
Мұндағы $V_\odot$ --- Күннің толық көлемі. \textit{Кеңес:} \taskref{7}-пунктте сізге $\int_a^b u\,dv= uv|_a^b-\int_a^b v\,du$ бөліктеп интегралдау әдісі қажет болуы мүмкін.
\begin{task}
\item Күннің толық гравитациялық энергиясы $E_G$ мен орташа қысым $\langle P\rangle$ келесі өрнекпен байланысқандығын көрсетіңіз:
$$
E_G = \beta\langle P\rangle V_\odot,
$$
мұндағы $\beta$ --- өлшемсіз сандық параметр. $\beta$ мәнін табыңыз.
\item Күннің тығыздығы орталыққа жақындаған сайын артатындықтан, $E_G$ өрнегіндегі $f$ коэффициенті біртекті бұлт үшін \taskref{4}-пункттегі жауаптан өзгеше болады. Күн үшін $f\approx1$ деп жорамалдап, Күн ішіндегі орташа қысым $\langle P\rangle$-дің сандық мәнін бағалаңыз.
\end{task}
Күн ішіндегі орташа $T_0$ температурасын бағалайық. Ол үшін $\langle P\rangle$ қысымы мен тығыздық Күн ішінде біртекті таралған деп және Күн затын электрондардан, иондалған сутегі мен гелийден тұратын идеал газ деп есептейік. Иондалған газ бөлшектерінің орташа массасы шамамен $\overline m=\qty{0.61}{\amu}$.
\begin{task}
\item Жоғарыдағы жорамалдарды ескере отырып, Күн ішіндегі $T_0$ температурасын сандық түрде бағалаңыз.
\end{task}

\part{Күн сәулеленуі}

Бұл есептің соңғы бөлімінде Күннің сәулелену қасиеттерін зерттейміз. Күннің, басқа жұлдыздардың көпшілігі сияқты, жылулық сәулелену спектрі абсолют қара дененің спектріне өте жақын екені, ал Күннің жарқырауы (яғни толық сәулелену қуаты) $L_\odot=\qty{4e26}{\W}$ екені белгілі.
\begin{task}
\item Күн бетіндегі $T_s$ температурасын бағалаңыз.
\end{task}

\pic{11}{R}{0pt}{0.27}{figures/2.4.png}

Бірінші жуықтауда Күнді $T_0$ температурасында радиациямен термодинамикалық тепе-теңдікте тұрған электрондар мен иондардан тұратын шар ретінде қарастыруға болады --- онда Күн $T_0$ температурасындағы абсолют қара радиатор болар еді және сәйкесінше басқа $L_\odot'$ жарқырауына ие болар еді. $L_\odot$ мен $L'_\odot$ арасындағы айырмашылық Күн ішіндегі радиацияның электрондар мен иондар тарапынан үнемі шашырауына, жұтылуына және қайта сәулеленуіне байланысты туындайды.

\textit{Радиативті диффузия} құбылысын модельдеу үшін, Күн орталығынан қозғала бастаған фотон жұтылу/сәулелену салдарынан сынық сызық бойымен жүреді деп елестетейік, мұндағы еркін жүру жолының орташа ұзындығы $l$-ді фотон үшін тұрақты деп санауға болады. Келесі жұтылу алдындағы кезекті $\vec l_i$ ығысу бағыты мүлдем кездейсоқ және алдыңғы $\vec l_{i-1}$ орын ауыстыруына тәуелді емес, сонымен қатар барлық $i$ үшін $|\vec l_i|=l$.

\begin{task}
\item Фотонның толық орын ауыстыру векторы $\vec D$-ны $\vec l_1,\dots,\vec l_N$ арқылы жазыңыз. Фотонның $N$ ығысу қадамынан кейінгі орташа орын ауыстыруы $\langle\vec D\rangle$ неге тең?
\item Орташа квадраттық орын ауыстыру $\langle D^2\rangle$-ны $N$ және $l$-ге тәуелділікте бағалаңыз. Осыдан фотонның Күннен шығуы үшін қажетті бағалау уақыты $t$-ны табыңыз.
\end{task}
Радиативті диффузияның арқасында біз $L_\odot \approx L_\odot'\cdot l/R_\odot$ қатынасын аламыз. Бұл байланыс бізге $L_\odot$ жарқырауының Күн массасы $M_\odot$-ға қалай тәуелді екенін түсінуге көмектеседі.
\begin{task}
\item $l$-дің сандық мәнін $T_0$, $T_s$ және $R_\odot$-ге тәуелділікте бағалаңыз. Осыдан \taskref{12}-пункттегі $t$ уақытының мәнін бағалаңыз. Нәтиже Күннің ішкі бөлігінің өте мөлдір емес екенін растайды.
\item \taskref{9}-пунктінің нәтижесін қолдана отырып, $L_\odot$ мен $M_\odot$ бір-бірімен келесі түрде байланысқанын көрсетіңіз:
$$
L_\odot = C\cdot\langle\rho\rangle\cdot(M_\odot)^n.
$$
Мұндағы $\langle\rho\rangle=\frac{M_\odot}{V_\odot}$. Сандық есептеулерсіз $C$ коэффициентінің $\overline m$, $l$ параметрлерімен, сондай-ақ іргелі тұрақтылармен байланысын көрсетіңіз және $n$ дәрежесінің сандық мәнін анықтаңыз.
\end{task}
\taskref{14}-пункті нәтижесінің маңыздылығын атап өткен жөн. $L_\odot$-ның $M_\odot$-ға дәрежелік тәуелділігі бас тізбектегі жұлдыздар үшін өте жақсы байқалады. $0.2M_\odot$--$50M_\odot$ аралығындағы жұлдыздар үшін масса--жарқырау тәуелділік графигі төменде көрсетілген ($\log$ негізі 10 бойынша алынған).
\begin{task}
\item График негізінде \taskref{14}-пунктінде анықталған $n$ дәрежесін қателікті есептемей бағалаңыз.
\end{task}

\cpic{1}{figures/2.5.png}[$0.2M_\odot$--$50M_\odot$ аралығындағы масса--жарқырау тәуелділігі. (Torres et al. (2010))]
